% Generated draft: requires mathematical and source review.
\documentclass[11pt]{article}
\usepackage[margin=25mm]{geometry}
\usepackage{fontspec}
\setmainfont{Latin Modern Roman}
\usepackage{amsmath,amssymb,mathtools,unicode-math}
\setmathfont{Latin Modern Math}
\usepackage{enumitem,needspace,xurl,hyperref}
\setlength{\parindent}{0pt}
\setlength{\parskip}{5pt}
\setlength{\emergencystretch}{3em}
\newcommand{\N}{\mathbb{N}}
\newcommand{\Z}{\mathbb{Z}}
\newcommand{\Q}{\mathbb{Q}}
\newcommand{\R}{\mathbb{R}}
\newcommand{\C}{\mathbb{C}}
\newcommand{\F}{\mathbb{F}}
\newcommand{\A}{\mathbb{A}}
\newcommand{\PP}{\mathbb{P}}
\newcommand{\E}{\mathbb{E}}
\newcommand{\eps}{\varepsilon}
\newcommand{\dd}{\,\mathrm{d}}
\begin{document}
\section*{UnsolvedMath ID: 701602}
\textbf{Source problem code:} catalogue entry
\subsection{Definitions and mathematical statement}
Let \(G\) be a finite group. A subgroup \(A<B\) is maximal in \(B\) if there is no subgroup \(L\) with \(A<L<B\). For an integer \(m\geq 1\), a subgroup \(H\leq G\) is called \(m\)-maximal in \(G\) if there exists a chain
\[
H=H_0<H_1<\cdots<H_m=G
\]
in which \(H_i\) is maximal in \(H_{i+1}\) for every \(0\leq i<m\). This definition requires the existence of a chain of length \(m\); a subgroup may satisfy it for more than one value of \(m\).

A chief factor of \(G\) is a quotient \(N/K\), where \(K,N\trianglelefteq G\), \(K<N\), and there is no normal subgroup \(L\trianglelefteq G\) satisfying \(K<L<N\).

Determine the chief factors of finite groups \(G\) satisfying
\[
\forall H\leq G,\qquad
\bigl(H\text{ is }2\text{-maximal in }G\bigr)
\Longrightarrow
\bigl(\forall m\in\Z,\ m\geq3,\ H\text{ is not }m\text{-maximal in }G\bigr).
\]
In particular, characterize the groups \(N/K\) that can occur as chief factors of such a \(G\). If \(G\) has no \(2\)-maximal subgroups, the displayed condition holds vacuously.

The source additionally states that for every integer \(m\geq3\), there exists a finite group \(G_m\) with a subgroup that is both \(2\)-maximal and \(m\)-maximal in \(G_m\). This is a supplied contextual assertion, not a conclusion proved here. The source poses a classification question and supplies no proposed characterization.
\subsection{Short English statement}
Determine the chief factors of a finite group in which every subgroup lying two maximal inclusions below the whole group admits no chain of maximal inclusions to the whole group of any length at least three. The source attributes this question to V. S. Monakhov.
\subsection{Sources}
\begin{itemize}
\item Supplied local UnsolvedMath registry record, permanent ID \(701602\), exact source problem code \texttt{catalogue entry}, titled \emph{Kourovka Notebook Problem 19.54}. The record labels the problem \emph{open}; this is supplied registry status, not an independently verified assessment.
\item V. S. Monakhov, Problem 19.54, \emph{The Kourovka Notebook (Unsolved Problems in Group Theory)}, 19th issue (2018). The supplied background also identifies the 21st issue edition (2026), \texttt{arXiv:1401.0300v48}. The registry describes the collection as edited by E. I. Khukhro and V. D. Mazurov. These bibliographic details and attribution are supplied and have not been independently checked.
\end{itemize}
\end{document}
